\documentclass{missingbook}
\usepackage{singular}

% chapters are included when they exist, so that the book builds while it is being written
\newcommand{\chapterfile}[1]{\IfFileExists{chapters/#1.tex}{\input{chapters/#1}}{}}

\begin{document}

\frontmatter
\hypersetup{pdftitle={Singular Learning Theory for the Rest of Us: From y = abx to Neural Networks},pdfauthor={Written by Claude; conceived and edited by Xiaoyu Li}}
\mbtitlepage{Singular Learning Theory for the Rest of Us}{From $y=abx$ to Neural Networks}{Written by Claude\\[2.5mm]Conceived and edited by Xiaoyu Li}{First edition \textperiodcentered{} October 2026}
\IfFileExists{frontmatter/preface.tex}{\addchap{Preface}

This book began with a list of textbooks that should exist but do not. One entry observed that neural networks are singular statistical models: many different weights give the same function, the Fisher information is degenerate, and the asymptotic theory behind the Laplace approximation, the Bayesian information criterion, and Akaike's criterion does not apply to them. The entry noted that the correct theory exists. Sumio Watanabe built it from 1999 on, and it replaces the familiar half the number of parameters by a single invariant of the loss, the learning coefficient. It has recently been used to measure how the complexity of a network changes during training. But Watanabe's monograph of 2009 asks its readers to know resolution of singularities before they begin, and the theory has a reputation for being hard to learn. The entry proposed a textbook for the rest of us: one family of toy models, starting from the model $y=abx$ with two parameters, on which every concept is computed by hand; Hironaka's theorem on resolution of singularities used as a black box, stated precisely and never proved; and a last part that meets the empirical results of deep learning.

The resulting book has fifteen chapters in four parts. Part I sets the classical theory against the running toy models and computes the free energy of the simplest of them completely, which turns out to contain most of the theory in miniature. Part II is about the learning coefficient as a geometric quantity: the rate at which the volume of nearly true parameters shrinks, computed by blowing up the parameter space or by reading a Newton polygon, with a catalogue of the values that are known. Part III is about Watanabe's theorems on the free energy, the generalization error, the criteria WAIC and WBIC, and the posterior's response to temperature: it proves the central result of each completely for a class of models where the computation can be carried through, states the general theorem with its source, and shows how a posterior moves between singularities as the data grow. Part IV asks what all this says about trained networks: the local learning coefficient, its changes during training, and what singular learning theory does not explain. Each chapter states what the regular theory says about its subject, computes the singular answer by hand on a toy model, proves its central theorem completely, in general or for such a class of models, measures the result in a short script, records it on a ledger card, and closes by saying where the theory stops.

The intended readers are researchers and graduate students in machine learning and statistics who have met the claim that neural networks are singular and want the theory behind it. The prerequisites are multivariable calculus, linear algebra, and probability at the level of a first graduate course; no algebraic geometry is assumed, and none beyond the statement of one theorem is used. Each chapter defines what it needs, and an appendix collects the shared tools, with proofs of those whose proofs are short.

Two cautions belong at the front. First, singular learning theory is a theory of exact Bayesian posteriors as the number of examples grows with the model fixed. Deep learning trains one network by stochastic gradient descent, often with more parameters than examples. The book tries to say exactly which results transfer, which transfer under assumptions that can be checked, and which do not transfer at all; the last chapter is devoted to the gap. Second, much of the work that connects the theory to deep learning is young. Several of the results discussed are preprints from 2023 to 2026, and most of the experiments behind them are small. The chapters date such claims and say when a result has not been refereed.

\paragraph{A note on authorship.} This book was written by Claude, an AI model made by Anthropic, and conceived and edited by Xiaoyu Li, who proposed the subject, its audience, the gap it fills, and its approach through one family of toy models. Every theorem is stated with its hypotheses and is proved, sketched, or attributed to the source that proves it; several results and proofs are the book's own and are marked as such. Every claim attributed to a paper was checked against the paper or a reliable account of it wherever these could be reached, every number from the book's scripts comes from code that the reader can run, and the colophon describes how the book was made. Readers who find an error are invited to report it, with the computation that fails.
}{}
\mbtableofcontents
\IfFileExists{frontmatter/notation.tex}{\addchap{Notation}

The chapters share the notation below. A chapter that needs a clashing local symbol says so where it introduces it, and the notes at the end of each chapter translate the notation of its main sources into this one.

\addsec*{Data, models and losses}

\begin{tabular}{@{}p{0.27\textwidth}p{0.69\textwidth}@{}}
$X^n=(X_1,\dots,X_n)$ & the data, independent with distribution $q$, the truth \\
$(x_i,y_i)$ & an input and an output in a regression $y=f(w,x)+\varepsilon$, with $\varepsilon\sim\Normal(0,\sigma^2)$ and $\sigma=1$ unless stated \\
$p(x\mid w)$, $w\in W\subset\R^d$ & a model and its parameter; $W$ is compact unless stated \\
$\varphi(w)$ & the prior density, uniform on $[-1,1]^d$ unless stated \\
$a,b$; $a_1,\dots,a_L$; $A,B$ & the parameters of the running models (two scalars; a chain; two matrices) \\
$L_n(w)$, $L(w)$ & the empirical loss $-\frac1n\sum_i\log p(X_i\mid w)$ and its expectation \\
$S$, $S_n$ & the entropy of $q$ and its empirical version $-\frac1n\sum_i\log q(X_i)$ \\
$w_0$, $W_0$ & a true (optimal) parameter, and the set of all of them \\
$K(w)$, $K_n(w)$ & $L(w)-L(w_0)\ge0$, the Kullback--Leibler divergence $\KL(q\,\|\,p(\cdot\mid w))$ in the realizable case, and $L_n(w)-L_n(w_0)$
\end{tabular}

\addsec*{Posteriors and free energies}

\begin{tabular}{@{}p{0.27\textwidth}p{0.69\textwidth}@{}}
$\beta$ & an inverse temperature; the tempered posterior has density proportional to $\varphi(w)e^{-n\beta L_n(w)}$ \\
$\Ew{f}$, $\Ewb{f}$ & the posterior expectation of $f(w)$ at $\beta=1$ and at $\beta$ \\
$p^*(x)$ & the Bayes predictive density $\Ew{p(x\mid w)}$ \\
$F_n$, $F_n(\beta)$ & the free energy $-\log\int\varphi(w)e^{-nL_n(w)}\dd w$ (the negative log marginal likelihood), and its tempered version \\
$F_n-nS_n$ & the normalized free energy $-\log\int\varphi(w)e^{-nK_n(w)}\dd w$ \\
$\bar F(n)$ & the population free energy $-\log\int\varphi(w)e^{-nK(w)}\dd w$
\end{tabular}

\addsec*{The learning coefficient}

\begin{tabular}{@{}p{0.27\textwidth}p{0.69\textwidth}@{}}
$\zeta(z)$ & the zeta function $\int K(w)^z\varphi(w)\dd w$ \\
$\lambda$, $m$ & the learning coefficient (the real log canonical threshold of $K$ with respect to $\varphi$) and its multiplicity: $-\lambda$ is the largest pole of $\zeta$ and $m$ its order \\
$V(\eps)$, $v(t)$ & the volume $\int_{K(w)<\eps}\varphi(w)\dd w$ and the state density $V'(t)$ \\
$\lambda(w^*)$, $\hat\lambda(w^*)$ & the local learning coefficient at a point $w^*$ and an estimate of it \\
$\nu$ & the singular fluctuation
\end{tabular}

\addsec*{Generalization and criteria}

\begin{tabular}{@{}p{0.27\textwidth}p{0.69\textwidth}@{}}
$G_n$, $T_n$ & the Bayes generalization loss $-\E_X[\log p^*(X)]$ and the Bayes training loss $-\frac1n\sum_i\log p^*(X_i)$; $G_n-S$ is the generalization error \\
$V_n$ & the functional variance $\sum_i\Var_w[\log p(X_i\mid w)]$ \\
$\WAIC$, $\WBIC$ & $T_n+V_n/n$, and $\Ewb{nL_n(w)}$ at $\beta=1/\log n$; both on the scale of the loss, never multiplied by $2n$ \\
$\AIC$, $\BIC$ & the classical criteria, written on the same scales
\end{tabular}

\addsec*{Probability and asymptotics}

$\Prob$ and $\E$ without a subscript are over the data (and fresh test points); a subscript names other randomness. $\Normal(\mu,\sigma^2)$ is the normal distribution and $\Unif$ the uniform one. $O_p$ and $o_p$ refer to $n\to\infty$ with the randomness of the data; $a_n\sim b_n$ means $a_n/b_n\to1$, and $a\asymp b$ means that each is bounded by a constant multiple of the other. Logarithms are natural. $\N$ is the set of natural numbers.
}{}

\mainmatter
\chapterfile{00-prologue}

\casepart{Where Regular Theory Ends}{Classical statistics charges a model $\tfrac d2\log n$ for its $d$ parameters, and it is right when the parameters are identifiable and the Fisher information is invertible. Neural networks satisfy neither condition. Two parameters are enough to see what goes wrong and to compute, by hand, what replaces the classical answer.}{{2/2}{2}{\tfrac12}{a^2b^2}}
\chapterfile{01-regular}
\chapterfile{02-singular}
\chapterfile{03-coefficient}

\casepart{The Learning Coefficient}{The number that replaces $\tfrac d2$ is the learning coefficient $\lambda$: the rate at which the volume of nearly true parameters shrinks. It is computed by changing coordinates until the loss becomes a monomial, which is what resolution of singularities guarantees and what Newton polygons do in one step.}{{0/4,2/0}{4/3}{\tfrac34}{a^2+b^4}}
\chapterfile{04-volume}
\chapterfile{05-blowups}
\chapterfile{06-newton}
\chapterfile{07-catalogue}

\casepart{Watanabe's Theorems}{With the learning coefficient in hand, the classical formulas can be rewritten for singular models: the free energy, the generalization error, the information criteria, and the response of the posterior to temperature. The posterior does not settle on one point; it moves between singularities as data arrive.}{{0/6,2/2,6/0}{2}{\tfrac12}{a^6+a^2b^2+b^6}}
\chapterfile{08-freeenergy}
\chapterfile{09-generalization}
\chapterfile{10-waic}
\chapterfile{11-temperature}
\chapterfile{12-phases}

\casepart{Deep Learning}{A trained network is not a posterior, and a large one is not in the asymptotic regime of the theorems. Still, the learning coefficient can be estimated at a trained network, and it moves during training. This part says what such measurements show, and what singular learning theory does not explain.}{{2/4,4/2}{3}{\tfrac13}{a^4b^2+a^2b^4}}
\chapterfile{13-llc}
\chapterfile{14-stages}
\chapterfile{15-sgd}

\bookmarksetup{startatroot}
\chapterfile{16-epilogue}

\appendix
\chapterfile{17-toolkit}

\backmatter
\printindex
\IfFileExists{backmatter/colophon.tex}{\specialchapter{How This Book Was Made}{Colophon}{{0/6,4/0}{12/5}{\tfrac5{12}}{a^4+b^6}}

\addsec{Roles}

\emph{Singular Learning Theory for the Rest of Us} was written by Claude Opus 5.5, an AI model made by Anthropic, in October 2026. Xiaoyu Li conceived the series \emph{The Missing Textbooks} and proposed this volume: its subject and its audience, researchers in deep learning; its gap (a theory now used to measure how neural networks change during training, without a textbook that such readers find easy to learn from); its approach (one family of toy models, starting from the two-parameter model $y=abx$, on which every concept is computed by hand, with Hironaka's theorem used as a black box and the last chapters meeting the empirical results of deep learning); and its position relative to existing books (Watanabe's two books, the exercise books on WAIC and WBIC, and a series of blog posts). Xiaoyu Li also set the direction of the series (language, design, and credit) and is responsible for the book's publication.

\addsec{Process}

The book was written in one long session at maximum reasoning effort. A coordinating instance of the model first surveyed the literature with four research assistants, separate instances that verified, with web search and the papers themselves, Watanabe's core theory, the computation of learning coefficients, the work that applies the theory to deep learning, and the classical theory together with the books that already exist; they found no book, lecture notes, or course that teaches the theory to deep-learning researchers by computing every quantity on a family of small models. Their digests corrected or sharpened the plan in several places before a chapter was written: with Gaussian inputs the loss of the tanh neuron is smooth but not analytic, so the book draws its inputs from a uniform law; the mixture's loss is not analytic in its weight; the upper half of the Hessian bounds on the learning coefficient had been published in 2011 and 2017, before it was rediscovered in work on deep learning; WBIC differs from the free energy by a term of order $\sqrt{\log n}$, not by a bounded one; and that the estimator of the local learning coefficient with a Gaussian localizer recovers it is a conjecture of the estimator's authors, not a theorem.

The coordinator then wrote the plan (the thesis, the running models, the notation, fixed labels for results that chapters cite across the book, and a specification for each chapter), a style guide (the anatomy of a chapter, the voice, a proof style that follows Halmos, and the conventions of the scripts), the volume's typography, a shared kit for the scripts, and the prologue.

Separate instances of the same model drafted the fifteen chapters and the appendix, up to thirteen at a time. Each verified its references and the claims it attributes to papers, wrote and ran the scripts behind its figures, checked its theorems and constants numerically, and built its chapter with the volume's checks for errors, references, layout, and style. Several authors found that their specifications were wrong, and said so. The correction term in the tempered mean energy of \cref{ch:temperature} does not follow from the asymptotics of the volume alone, as a slowly oscillating counterexample shows, so the chapter adds an assumption that analytic losses satisfy. The two regions of the mixture specified for \cref{ch:phases} do not separate, and the chapter replaced the mixture with a double well. The digest's account of Watanabe's non-renormalizable example was wrong: the source states that the equations of state still hold, at order $n^{-2/3}$ (\cref{ch:waic}), and the author of \cref{ch:waic} and the referees of \cref{ch:freeenergy,ch:waic} found that two constants printed in the source are too large by the factor $2^{4/3}$. Several results are, as far as the authors could find, new; the plan proposed some of them, which the authors proved and in several cases strengthened, and the authors found the others. Among them are the law of the posterior of $y=abx$ over scales (\cref{ch:singular}); the exact generalization error $1/(n-1)$ of the regular plane under a flat prior (\cref{ch:regular}); a proof that $n$ times the expected generalization error tends to $\lambda$ for every one-coefficient regression with Gaussian inputs and noise, through Tweedie's formula, a virial identity, and Stein's identity (\cref{ch:generalization}); a proof that leave-one-out cross-validation and WAIC agree to order $n^{-2}$ for $y=abx$ (\cref{ch:waic}); the tempered mean energy with its correction term and a complete proof of the WBIC formula for realizable linear models (\cref{ch:temperature}); a proof of the local free energy formula for models linear in their coefficients, in regions that need not contain the truth (\cref{ch:phases}); a proof that with a fixed Gaussian localizer the population version of the estimator of the local learning coefficient can converge to the coefficient of another point (\cref{ch:llc}); a form of the rule that a minimally singular set of critical points has learning coefficient half its codimension, which derives that the set is a minimum instead of assuming it, with a Fourier computation of the Hessian at the hexagon of the toy model of superposition and a calibration of the estimator at its polygons (\cref{ch:stages}); and an analysis of where gradient flow on $y=abx$ stops, which adds convergence and a scale to the known conservation law (\cref{ch:sgd}). The chapters mark these as the book's own.

Each chapter, and the prologue and appendix together, then received one independent referee pass from a separate instance, which re-derived the results line by line and checked them numerically, reran the scripts and compared every printed number with the text, worked the exercises, and checked the claims about papers against the papers; it corrected what it found directly. No referee found a main theorem that was false as stated, but the referees found real errors. Some hypotheses were wrong: a lemma on importance sampling lacked two conditions without which its identity can fail (\cref{ch:waic}), a theorem on Newton polyhedra did not require the prior to sit near the point (\cref{ch:newton}), a rule for products of losses omitted a convergence hypothesis (\cref{ch:blowups}), and a lemma on symmetric criticality assumed a differentiability that the toy model of superposition, to which it was applied, does not have (\cref{ch:stages}). A proof about a double well had a gap (\cref{ch:phases}). The transition map between two charts of a blow-up was wrong, and the second term in the volume of the cusp had the wrong sign and size (\cref{ch:blowups}). A worked example claimed that Hessian ranks cannot show which of two points is more singular, which is false (\cref{ch:llc}). Some claims outran their evidence: one chapter said that it proves Watanabe's equations of state for one-coefficient models, where it proves three of the four limits they relate, and that it confirms the learning coefficient of a model whose measured interval excludes it (\cref{ch:generalization}); three said that sets of estimates lay inside their intervals when some or all lie outside (\cref{ch:temperature,ch:llc,ch:stages}); and an error estimate presented as a bound was not a bound, and missed the typical error by five orders of magnitude (\cref{ch:regular}). In \cref{ch:catalogue} a validation run contradicted a value proved in the same chapter, through the bias of a random-walk sampler, and a measured excess was attributed to the wrong cause; in \cref{ch:stages} the calibration intervals ignored the variation between sampler seeds; and a table caption in \cref{ch:volume} claimed checks that no script performed, until the referee wrote them. A heuristic for the collapse of stochastic gradient descent onto a saddle claimed more than holds off an invariant line (\cref{ch:sgd}). Several exercises could not be solved as posed, and two were false. The referees of \cref{ch:phases,ch:stages} confirmed that a crossing point reported for the toy model of superposition charges the loss twice what the paper's own likelihood does, which moves it from about $600$ samples to about $1300$; both chapters report it as an apparent inconsistency in a preprint. Two claims of novelty were too broad: the referee of \cref{ch:phases} found work from 2003 on the scaling that the chapter had presented as new, and the referee of \cref{ch:sgd} found one of the chapter's identities in a 2018 paper; both claims are now qualified. Each error is corrected in the text. No referee found an arithmetic bug in a script, but many quoted numbers were misrounded or out of date with their scripts.

A usage limit stopped the running agents once, authors and referees alike, and each was resumed with its context. Once the chapters were drafted, the coordinator wrote the preface, the notation, and the epilogue. It then integrated the book: it reconciled cross-references between chapters, gave every work cited in several chapters a single bibliographic entry, rebuilt the epilogue's ledger from the chapters' final cards, and checked the whole book for layout, undefined references, and overflowing lines. A last referee pass checked the epilogue, the preface, the notation, this colophon, and the volume's README against the chapters.

\addsec{Limits}

No human has yet checked every proof or every attribution. Web search and the papers themselves were available while the book was written, but some full texts could not be opened: Watanabe's 2009 book is cited by chapter and section, and by theorem number only as later works, his own papers among them, cite it, and some claims about what a paper proves were checked against abstracts, the paper's own summary, or secondary accounts; the chapters say so where it matters. One referee pass per chapter finds many errors but not all of them, and a referee that is an instance of the same model may share the author's blind spots. Results that the book presents as its own have been checked by the same model only. Much of the work that connects the theory to deep learning consists of preprints from 2023 to 2026, and most of the experiments behind it are small. The book's scripts are small too: most compute exact posteriors of models with a handful of parameters, and the rest run samplers or gradient descent on models with at most a few dozen, which confirms the arithmetic of the theorems, not their relevance to large networks. Corrections are welcome and will be recorded in the repository of \emph{The Missing Textbooks}.

\medskip
{\small\noindent Typeset with pdf\LaTeX{} in \texttt{missingbook}, the series class built on \textsc{koma}-Script, with the volume style \texttt{singular}; Palatino (newpx) for text and mathematics, Source Sans for labels, biblatex and biber for references, and Ti\emph{k}Z and pgfplots for figures, whose data were computed with Python, NumPy, and SciPy, and checked with mpmath and SymPy. The title page carries $2400$ draws from the posterior of $y=abx$ at its singular truth, computed on a grid of $1200\times1200$ cells, for $400$ observations of pure noise (among the data sets with seeds $0$ to $99$, the one whose least-squares slope is closest to $0$): the cross of \cref{fig:prologue:draws}. The foot of each chapter's opening page, and of each part page, carries the Newton polygon of a loss, the diagonal, and the point $(t,t)$ at which the diagonal meets the polygon, which gives the learning coefficient $1/t$ (\cref{ch:newton}); this page carries that of the loss $a^4+b^6$, whose coefficient is $\frac5{12}$. Both are computed or drawn by \texttt{figures/scripts/00-emblem.py} and \texttt{singular.sty}.\par}
}{}
\end{document}
